\documentclass[11pt,a4paper]{article}
\newcommand{\intestazione}{Dimero di spin --- 9. Scan sui parametri di rumore}
% =====================================================================
%  Preambolo comune ai documenti sul dimero (serie "che cosa abbiamo fatto")
%  Incluso con % =====================================================================
%  Preambolo comune ai documenti sul dimero (serie "che cosa abbiamo fatto")
%  Incluso con % =====================================================================
%  Preambolo comune ai documenti sul dimero (serie "che cosa abbiamo fatto")
%  Incluso con \input{preambolo_dimero} subito dopo \documentclass.
% =====================================================================
\usepackage[T1]{fontenc}
\usepackage[utf8]{inputenc}
\usepackage[main=italian,provide=*]{babel}
\usepackage{amsmath,amssymb,mathtools}
\usepackage{graphicx}
\usepackage{booktabs}
\usepackage{colortbl}
\usepackage{xcolor}
\usepackage{geometry}
\usepackage{placeins}
\usepackage{caption}
\usepackage{enumitem}
\usepackage{fancyhdr}
\usepackage{needspace}
\usepackage{tikz}
\usetikzlibrary{arrows.meta,positioning,shapes.geometric,calc,fit,backgrounds}
\usepackage{hyperref}
\hypersetup{colorlinks=true,linkcolor=blu,citecolor=blu,urlcolor=blu}

\geometry{margin=2.5cm}

% --- colori coerenti con quelli delle figure (palette Okabe-Ito)
\definecolor{blu}{HTML}{0072B2}
\definecolor{arancio}{HTML}{E69F00}
\definecolor{verdes}{HTML}{009E73}
\definecolor{rossos}{HTML}{D55E00}
\definecolor{violas}{HTML}{CC79A7}
\definecolor{grigetto}{HTML}{E8E8E8}

\captionsetup{font=small,labelfont=bf,width=0.92\textwidth}

\pagestyle{fancy}
\fancyhf{}
\fancyhead[L]{\small\itshape\intestazione}
\fancyhead[R]{\small\thepage}
\renewcommand{\headrulewidth}{0.3pt}

% --- notazione
\newcommand{\ket}[1]{\lvert #1\rangle}
\newcommand{\bra}[1]{\langle #1 \rvert}
\newcommand{\media}[1]{\langle #1 \rangle}
\newcommand{\vs}[1]{\mathbf{s}_{#1}}

% --- riquadro "in una riga": il messaggio da portare a casa
\newcommand{\messaggio}[1]{%
  \begin{center}
  \begin{minipage}{0.93\textwidth}
  \setlength{\fboxsep}{9pt}%
  \colorbox{grigetto}{\begin{minipage}{\dimexpr\textwidth-2\fboxsep}
  \small\textbf{In una riga.}\enspace #1
  \end{minipage}}
  \end{minipage}
  \end{center}
}

% --- riquadro "come lo sappiamo": la verifica numerica dietro un'affermazione
\newcommand{\verifica}[1]{%
  \begin{center}
  \begin{minipage}{0.93\textwidth}
  \small\textcolor{verdes}{\rule{2.2em}{1.1pt}}\enspace
  \textcolor{verdes}{\textbf{Come lo sappiamo.}}\enspace #1
  \end{minipage}
  \end{center}
  \medskip
}
 subito dopo \documentclass.
% =====================================================================
\usepackage[T1]{fontenc}
\usepackage[utf8]{inputenc}
\usepackage[main=italian,provide=*]{babel}
\usepackage{amsmath,amssymb,mathtools}
\usepackage{graphicx}
\usepackage{booktabs}
\usepackage{colortbl}
\usepackage{xcolor}
\usepackage{geometry}
\usepackage{placeins}
\usepackage{caption}
\usepackage{enumitem}
\usepackage{fancyhdr}
\usepackage{needspace}
\usepackage{tikz}
\usetikzlibrary{arrows.meta,positioning,shapes.geometric,calc,fit,backgrounds}
\usepackage{hyperref}
\hypersetup{colorlinks=true,linkcolor=blu,citecolor=blu,urlcolor=blu}

\geometry{margin=2.5cm}

% --- colori coerenti con quelli delle figure (palette Okabe-Ito)
\definecolor{blu}{HTML}{0072B2}
\definecolor{arancio}{HTML}{E69F00}
\definecolor{verdes}{HTML}{009E73}
\definecolor{rossos}{HTML}{D55E00}
\definecolor{violas}{HTML}{CC79A7}
\definecolor{grigetto}{HTML}{E8E8E8}

\captionsetup{font=small,labelfont=bf,width=0.92\textwidth}

\pagestyle{fancy}
\fancyhf{}
\fancyhead[L]{\small\itshape\intestazione}
\fancyhead[R]{\small\thepage}
\renewcommand{\headrulewidth}{0.3pt}

% --- notazione
\newcommand{\ket}[1]{\lvert #1\rangle}
\newcommand{\bra}[1]{\langle #1 \rvert}
\newcommand{\media}[1]{\langle #1 \rangle}
\newcommand{\vs}[1]{\mathbf{s}_{#1}}

% --- riquadro "in una riga": il messaggio da portare a casa
\newcommand{\messaggio}[1]{%
  \begin{center}
  \begin{minipage}{0.93\textwidth}
  \setlength{\fboxsep}{9pt}%
  \colorbox{grigetto}{\begin{minipage}{\dimexpr\textwidth-2\fboxsep}
  \small\textbf{In una riga.}\enspace #1
  \end{minipage}}
  \end{minipage}
  \end{center}
}

% --- riquadro "come lo sappiamo": la verifica numerica dietro un'affermazione
\newcommand{\verifica}[1]{%
  \begin{center}
  \begin{minipage}{0.93\textwidth}
  \small\textcolor{verdes}{\rule{2.2em}{1.1pt}}\enspace
  \textcolor{verdes}{\textbf{Come lo sappiamo.}}\enspace #1
  \end{minipage}
  \end{center}
  \medskip
}
 subito dopo \documentclass.
% =====================================================================
\usepackage[T1]{fontenc}
\usepackage[utf8]{inputenc}
\usepackage[main=italian,provide=*]{babel}
\usepackage{amsmath,amssymb,mathtools}
\usepackage{graphicx}
\usepackage{booktabs}
\usepackage{colortbl}
\usepackage{xcolor}
\usepackage{geometry}
\usepackage{placeins}
\usepackage{caption}
\usepackage{enumitem}
\usepackage{fancyhdr}
\usepackage{needspace}
\usepackage{tikz}
\usetikzlibrary{arrows.meta,positioning,shapes.geometric,calc,fit,backgrounds}
\usepackage{hyperref}
\hypersetup{colorlinks=true,linkcolor=blu,citecolor=blu,urlcolor=blu}

\geometry{margin=2.5cm}

% --- colori coerenti con quelli delle figure (palette Okabe-Ito)
\definecolor{blu}{HTML}{0072B2}
\definecolor{arancio}{HTML}{E69F00}
\definecolor{verdes}{HTML}{009E73}
\definecolor{rossos}{HTML}{D55E00}
\definecolor{violas}{HTML}{CC79A7}
\definecolor{grigetto}{HTML}{E8E8E8}

\captionsetup{font=small,labelfont=bf,width=0.92\textwidth}

\pagestyle{fancy}
\fancyhf{}
\fancyhead[L]{\small\itshape\intestazione}
\fancyhead[R]{\small\thepage}
\renewcommand{\headrulewidth}{0.3pt}

% --- notazione
\newcommand{\ket}[1]{\lvert #1\rangle}
\newcommand{\bra}[1]{\langle #1 \rvert}
\newcommand{\media}[1]{\langle #1 \rangle}
\newcommand{\vs}[1]{\mathbf{s}_{#1}}

% --- riquadro "in una riga": il messaggio da portare a casa
\newcommand{\messaggio}[1]{%
  \begin{center}
  \begin{minipage}{0.93\textwidth}
  \setlength{\fboxsep}{9pt}%
  \colorbox{grigetto}{\begin{minipage}{\dimexpr\textwidth-2\fboxsep}
  \small\textbf{In una riga.}\enspace #1
  \end{minipage}}
  \end{minipage}
  \end{center}
}

% --- riquadro "come lo sappiamo": la verifica numerica dietro un'affermazione
\newcommand{\verifica}[1]{%
  \begin{center}
  \begin{minipage}{0.93\textwidth}
  \small\textcolor{verdes}{\rule{2.2em}{1.1pt}}\enspace
  \textcolor{verdes}{\textbf{Come lo sappiamo.}}\enspace #1
  \end{minipage}
  \end{center}
  \medskip
}


\title{\bfseries Un solo chip non basta a dire qualcosa di generale\\[2pt]
\large Documento 9 --- come si sposta $N^*$ al variare del canale}
\author{Samuele Schiavi \\ \small Tesi triennale in Fisica, Università di Parma
--- Relatore: Prof.\ Alessandro Chiesa}
\date{}

\begin{document}
\maketitle
\thispagestyle{fancy}

\section{Perché non fermarsi al punto di calibrazione}

I due documenti precedenti hanno trovato $N^*$ a un \textbf{singolo} punto
della calibrazione di riferimento, \texttt{ibm\_torino}. Non è una domanda
oziosa chiedersi cosa succeda altrove: un chip diverso, o lo stesso chip in un
momento diverso, avrà calibrazioni diverse. Sapere come $N^*$ dipenda dai
parametri del canale di rumore rende il risultato utile oltre il singolo punto
già misurato --- ed è la prima volta, in questa serie, che si esplora un intero
spazio di parametri invece di un punto fisso.

\section{Una previsione, prima di guardare i numeri}

Il rumore di gate entra nella fedeltà o nel correlatore in modo non banale:
più forte è $\varepsilon_{1q}$ o $\varepsilon_{2q}$, più in fretta il rumore
accumulato supera il guadagno dell'errore di Trotter al crescere di $N$. Ci si
aspetta quindi che $N^*$ \textbf{diminuisca} al crescere del rumore di gate.

Il readout, invece, entra --- come visto nel Documento 8 --- come un fattore
puramente moltiplicativo, $(1-2p_\text{readout})$, applicato allo stesso modo a
tutti gli $N$. Se una funzione $g(N)$ ha un massimo a $N^*$, allora $c\cdot
g(N)$, per qualunque costante positiva $c$, ha il massimo allo \textbf{stesso}
$N^*$: moltiplicare per una costante positiva non può mai spostare un argmax.
La previsione, scritta prima di guardare i dati: $N^*$ dipende da
$\varepsilon_{1q}$ e $\varepsilon_{2q}$, non da $p_\text{readout}$.

\section{Monotonia rispetto al rumore di gate}

\begin{figure}[htbp]
\centering
\includegraphics[width=0.95\textwidth]{figure/fig_passo5_Nstar_vs_eps}
\caption{$N^*$ in funzione di $\varepsilon_{2q}$ (sinistra) e $\varepsilon_{1q}$
(destra), scala log, l'altro parametro fissato al valore di riferimento.
Linea tratteggiata: valore di calibrazione \texttt{ibm\_torino}. Monotonia
non-crescente confermata in entrambi i casi, con una transizione a $N^*=1$
per rumore sufficientemente forte.}
\end{figure}

\FloatBarrier

\begin{center}
\renewcommand{\arraystretch}{1.15}
\small
\begin{tabular}{@{}cc@{}}
\toprule
$\varepsilon_{2q}$ & $N^*$\\
\midrule
$1.0\times10^{-3}$ & $10$\\
$2.0\times10^{-3}$ & $9$\\
$3.8\times10^{-3}$ (riferimento) & $8$\\
$6.0\times10^{-3}$ & $7$\\
$1.0\times10^{-2}$ & $1$\\
$1.5\times10^{-2}$ & $1$\\
$2.5\times10^{-2}$ & $1$\\
\bottomrule
\end{tabular}
\qquad
\begin{tabular}{@{}cc@{}}
\toprule
$\varepsilon_{1q}$ & $N^*$\\
\midrule
$1.0\times10^{-4}$ & $9$\\
$2.9\times10^{-4}$ (riferimento) & $8$\\
$1.0\times10^{-3}$ & $7$\\
$3.0\times10^{-3}$ & $1$\\
$1.0\times10^{-2}$ & $1$\\
$3.0\times10^{-2}$ & $1$\\
\bottomrule
\end{tabular}
\end{center}

$N^*$ è monotono non-crescente in entrambi $\varepsilon_{1q}$ ed
$\varepsilon_{2q}$, come previsto. La transizione a $N^*=1$ non è un
artefatto numerico: oltre una certa soglia di rumore di gate, nessun numero di
passi di Trotter maggiore di 1 vale il costo aggiuntivo in gate rumorosi ---
un solo passo, per quanto impreciso, resta comunque la scelta meno rumorosa.

\section{Invarianza rispetto al readout}

\begin{center}
\renewcommand{\arraystretch}{1.2}
\begin{tabular}{@{}ccc@{}}
\toprule
$p_\text{readout}$ & $N^*$ & $|C(N^*)|$\\
\midrule
$0.00$ & $2$ & $0.6155$\\
$0.05$ & $2$ & $0.5534$\\
$0.20$ & $2$ & $0.3688$\\
\bottomrule
\end{tabular}
\end{center}

$N^*=2$ per tutti i valori di $p_\text{readout}$ testati, dallo 0\% al 20\% ---
un intervallo enormemente più ampio del valore di calibrazione reale (2.3\%).
Il valore assoluto $|C(N^*)|$ cambia in modo sensibile, da $0.62$ a $0.37$, ma
la \emph{posizione} del massimo no. A differenza di una prima versione di
questo documento, qui non è una conseguenza garantita dalla formula usata
per calcolare i numeri: ogni valore di questa tabella viene da una vera
istruzione di misura sull'ancilla, con un \texttt{ReadoutError} genuino
agganciato al simulatore e campionato a $2\times10^5$ shot (Documento 8) ---
la formula $(1-2p)\langle Z\rangle_\text{ideale}$ resta solo l'argomento
teorico che spiega \emph{perché} aspettarsi l'invarianza, ma questa tabella
la verifica su dati genuinamente campionati, non la presuppone.

\messaggio{La previsione teorica regge, verificata su dati genuinamente
campionati, non solo argomentata per plausibilità né garantita dalla
struttura del calcolo. Questa proprietà --- il readout simmetrico è sempre
un fattore moltiplicativo, mai additivo --- smette di valere non appena il
readout diventa asimmetrico, come si vedrà nel Documento 11: lì
l'argomento di invarianza usato qui non si estende automaticamente, ed è
per questo che quella domanda resta aperta fino a essere verificata
numericamente.}

\section{Un affinamento: quanto è rapida la saturazione a $N^*=1$?}

Un modello ingenuo per l'errore totale, $\varepsilon_\text{tot}(N)\sim
a/N+g\,\varepsilon_{2q}N$ (Trotter decrescente più rumore crescente, entrambi
additivi), predice per derivazione diretta un ottimo
$N^*=\sqrt{a/(g\varepsilon_{2q})}\propto\varepsilon_{2q}^{-1/2}$: una pendenza
$-0.5$ in scala log-log. Confrontando con i punti non ancora saturati della
tabella sopra ($\varepsilon_{2q}$ da $10^{-3}$ a $6\times10^{-3}$, $N^*$ da
$10$ a $7$):

\begin{center}
\renewcommand{\arraystretch}{1.2}
\begin{tabular}{@{}lc@{}}
\toprule
Quantità & Valore\\
\midrule
Pendenza log-log prevista (errore additivo) & $-0.5$\\
Pendenza log-log osservata (fit sui 4 punti non saturati) & $-0.195$\\
\bottomrule
\end{tabular}
\end{center}

Lo scostamento non indica un errore: indica che il modello additivo era
pensato per un errore su una singola quantità, mentre qui l'osservabile è una
\emph{fedeltà}. Con rumore depolarizzante accumulato su molti passi, una
fedeltà si comporta più come un decadimento moltiplicativo,
$F(N)\sim(1-\varepsilon_{2q})^{cN}$ per qualche costante $c$, che come una
somma di due termini indipendenti. La forma funzionale esatta di
$N^*(\varepsilon_{2q})$ per una fedeltà resta un affinamento possibile, non
necessario per rispondere alla domanda del relatore --- ``come si sposta
$N^*$'' --- a cui i dati sopra rispondono già direttamente.

\section{Cosa non copre questo scan}

Un parametro alla volta, con gli altri fissati al valore di riferimento: non
l'intero spazio congiunto $(\varepsilon_{1q},\varepsilon_{2q},p_\text{readout})$.
Non copre inoltre punti di lavoro fisici diversi da quello usato in tutta la
pipeline ($b/J=0.35$, $D/J=0.80$): se e quanto $N^*$ dipenda anche dal punto
fisico è una domanda verificata separatamente, con due punti di lavoro,
nell'estensione sul VQE noise-aware (Documento 10).

\section{Dove chiude questa parte della pipeline}

Cinque documenti, dal modello di rumore (due canali, quattro numeri, Documento
5) fino a qui, ognuno aggiungendo un livello di struttura: energia statica,
poi dinamica, poi un osservabile a due tempi che richiede una misura reale,
poi la sensibilità del risultato ai parametri stessi del rumore. Il filo
conduttore verificato a ogni stadio è lo stesso: il readout simmetrico è
sempre un fattore moltiplicativo, mai additivo --- proprietà che smette di
valere non appena il readout diventa asimmetrico (Documento 11), e che è alla
base di ogni argomento di invarianza usato finora.

\messaggio{$N^*$ dipende in modo monotono e prevedibile dal rumore di gate, e
non dipende affatto dal readout simmetrico --- entrambi i fatti verificati,
non solo previsti. I due documenti che seguono sono estensioni facoltative:
cosa succede rioptimizzando il VQE già sotto rumore (Documento 10), e cosa
succede quando il readout smette di essere simmetrico (Documento 11).}

\end{document}
